\chapter{Conception du système ESYAir}
% Ce chapitre est consacré à la description détaillée de l'implémentation de votre projet. Il a pour objectif d'expliquer, de manière claire et méthodique, comment vous avez concrètement mis en œuvre votre travail. Vous y présenterez, étape par étape, le processus qui a conduit à la réalisation de votre projet, en justifiant les choix techniques et méthodologiques effectués.

\index{expérience}
% Vous devez décrire le déroulement de vos expériences, en précisant le cadre dans lequel elles ont été menées ainsi que les outils et technologies utilisés. Ce chapitre doit permettre au lecteur de comprendre non seulement ce que vous avez fait, mais aussi pourquoi vous avez opté pour cette approche.

% Exposez de manière structurée la conception globale de votre projet en illustrant vos explications par des schémas, des diagrammes ou des modèles visuels -- tels que des diagrammes de classes, de flux ou des organigrammes -- afin de clarifier votre démarche.

% Que ce soit pour une architecture logicielle, une conception électronique ou une analyse de données, chaque aspect de votre travail doit être présenté de façon logique et détaillée. Il est essentiel de justifier vos choix en soulignant leur pertinence au regard des objectifs fixés et de la problématique posée.

% Vous êtes libre de structurer ce chapitre en plusieurs sections distinctes, en fonction de la nature de votre projet. L'objectif est d'offrir au lecteur une vision claire de la manière dont vous avez développé votre travail, en mettant en lumière non seulement les actions entreprises, mais également les décisions qui les ont motivées. Votre implémentation doit refléter une réflexion méthodique et rigoureuse, guidée par la problématique et les objectifs initiaux.
%------------------------------------------------------------------



% Ce chapitre présente la conception et l’architecture du système ESYAir.
% Il décrit les choix techniques et méthodologiques réalisés afin de répondre
% aux objectifs du projet. Les différentes composantes du système sont
% présentées ainsi que leur rôle dans la chaîne de mesure de la qualité de l’air.

\section{Architecture globale du système}

Description de l'architecture générale du système ESYAir comprenant
une station de référence au sol et un système mobile de mesure
permettant l'analyse verticale de la qualité de l'air.

\section{Station de référence au sol}

\subsection{Critères de sélection de la station}

Analyse multicritère utilisée pour sélectionner la station.

\subsection{Capteurs de gaz}

\subsection{Capteurs de particules}

\subsection{Capteurs météorologiques}

\section{Système mobile de mesure}

\subsection{Choix du vecteur de mobilité verticale}

Comparaison des solutions possibles :
\begin{itemize}
\item ballon météorologique
\item drone
\item autres solutions
\end{itemize}

\subsection{Capteurs embarqués}

\subsection{Système de positionnement GPS}

\section{Architecture de la chaîne de mesure}

Description du flux de données :
\begin{itemize}
\item acquisition
\item stockage
\item transmission
\end{itemize}